\newcounter{RDW} \setcounter{RDW}{1}
\newcounter{RDR} \setcounter{RDR}{1}
\newcounter{RDE} \setcounter{RDE}{1}
\newcounter{RDS} \setcounter{RDS}{1}
\newcounter{RF} \setcounter{RF}{1}
% Descripción del subsistema
\section*{Descripción}
El subsistema de Mensajería Privada se encarga de la comunicación privada entre usuarios, consistiendo principalmente en el envío, la eliminación y la visualización de mensajes que se dan entre dos usuarios del sistema, limitandose sólo a aquellos que son amigos. Se incluye la posibilidad de archivar chats o destacar mensajes para personalizar la visualización de los mensajes por parte del usuario. Parte desarrollada por \textbf{Sergio Calvo González}.
\newpage

% ======================= %
%     ENVIAR MENSAJE      %
% ======================= %
\section{Enviar mensaje}
\textbf{\textcolor{rojo}{\underline{RF5.\theRF:}}}. Enviar mensaje: Permite que un usuario envíe un mensaje de texto a otro usuario (deben ser amigos)\\[6pt]
\stepcounter{RF}

\textbf{\textcolor{rojo}{Entrada:}} Agente externo: usuario. Acción: solicitar envío de un mensaje. 
Requisito de datos de entrada \textcolor{azul}{RDE5.\theRDE}. \\[6pt]

\textbf{\textcolor{rojo}{BD:}} Requisito de datos de escritura \textcolor{azul}{RDW5.\theRDW}. \\[6pt]

\textbf{\textcolor{rojo}{Salida:}} Agente externo: usuario. Acción: confirmación resultado. 
Requisito de datos de salida: ninguno. \\[12pt]

\textcolor{azul}{\textbf{RDE5.\theRDE:}} Mensaje y usuario al que va dirigido \\
\hspace*{1cm} Usuario: Cadena de caracteres (20) \\
\hspace*{1cm} Mensaje: Cadena de caracteres (50) \\[12pt]
\stepcounter{RDE}

\textcolor{azul}{\textbf{RDW5.\theRDW:}} Datos del envío de mensaje. \\
\hspace*{1cm} Usuario Emisor: Cadena de caracteres (20) \\
\hspace*{1cm} Usuario Receptor: Cadena de caracteres (20) \\
\hspace*{1cm} Mensaje: Cadena de caracteres (50) \\[12pt]
\stepcounter{RDW}
% ======================= %
%    ELIMINAR MENSAJE     %
% ======================= %
\section{Eliminar mensaje}
\textbf{\textcolor{rojo}{\underline{RF\theRF:}}}. Eliminar mensaje: Permite que un usuario elimine un mensaje de texto que ha enviado, de forma que ni dicho usuario ni el receptor pueden volver a acceder al mensaje\\[6pt]
\stepcounter{RF}

\textbf{\textcolor{rojo}{Entrada:}} Agente externo: usuario. Acción: solicitar la eliminación de un mensaje. 
Requisito de datos de entrada: ninguno. \\[6pt]

\textbf{\textcolor{rojo}{BD:}} Requisito de datos de lectura \textcolor{azul}{RDR6.1}. \\[6pt]

\textbf{\textcolor{rojo}{Salida:}} Agente externo: usuario. Acción: confirmación resultado. 
Requisito de datos de salida: ninguno. \\[12pt]

\textcolor{azul}{\textbf{RDE6.2:}} Datos del mensaje a borrar\\
\hspace*{1cm} Usuario: Cadena de caracteres (20) \\
\hspace*{1cm} Mensaje: Cadena de caracteres (50) \\[12pt]

\textcolor{azul}{\textbf{RDR6.1:}} Lista de mensajes. \\
\hspace*{1cm} Usuario Emisor: Cadena de caracteres (20) \\
\hspace*{1cm} Usuario Receptor: Cadena de caracteres (20) \\
\hspace*{1cm} Mensaje: Cadena de caracteres (50) \\[12pt]

% ======================= %
%  LISTAR CONVERSACIONES  %
% ======================= %
\section{Listar conversaciones}
\textbf{\textcolor{rojo}{\underline{RF\theRF:}}}. Listar conversaciones: Permite que un usuario vea los conversaciones que tiene disponibles\\[6pt]
\stepcounter{RF}
\textbf{\textcolor{rojo}{Entrada:}} Agente externo: usuario. Acción: solicitar la visualización de sus conversaciones. 
Requisito de datos de entrada: ninguno\\[6pt]

\textbf{\textcolor{rojo}{BD:}} Requisito de datos de lectura \textcolor{azul}{RDR5.\theRDR}. \\[6pt]

\textbf{\textcolor{rojo}{Salida:}} Agente externo: usuario. Acción: confirmación resultado. 
Requisito de datos de salida: \textcolor{azul}{RDS5.\theRDS} \\[12pt]

\textcolor{azul}{\textbf{RDR5.\theRDR:}} Lista de usuarios con los que dispone de una conversación, con el formato:\\
\hspace*{1cm} Usuario: Cadena de caracteres (20) \\[12pt]
\stepcounter{RDR}

\textcolor{azul}{\textbf{RDS5.\theRDS:}} Lista de usuarios con los que dispone de una conversación. \\
\hspace*{1cm} Los mismos datos que (20) \\[12pt]
\stepcounter{RDS}
% ======================= %
%     VISUALIZAR CHAT     %
% ======================= %
\section{Visualizar chat}
\textbf{\textcolor{rojo}{\underline{RF\theRF:}}}. Visualizar chat: Permite que un usuario vea los mensajes entre otro usuario y él mismo\\[6pt]
\stepcounter{RF}
\textbf{\textcolor{rojo}{Entrada:}} Agente externo: usuario. Acción: solicitar la visualización de un chat. 
Requisito de datos de entrada \textcolor{azul}{RDE5.\theRDE}. \\[6pt]

\textbf{\textcolor{rojo}{BD:}} Requisito de datos de lectura \textcolor{azul}{RDR5.\theRDR}. \\[6pt]

\textbf{\textcolor{rojo}{Salida:}} Agente externo: usuario. Acción: confirmación resultado. 
Requisito de datos de salida: \textcolor{azul}{RDS6.2} \\[12pt]

\textcolor{azul}{\textbf{RDE5.\theRDE:}} Datos del usuario del chat que se desea ver\\
\hspace*{1cm} Usuario: Cadena de caracteres (20) \\[12pt]
\stepcounter{RDE}

\textcolor{azul}{\textbf{RDR5.\theRDR:}} Lista de mensajes entre ambos usuarios, ordenados temporalmente, con el siguiente formato:\\
\hspace*{1cm} Usuario emisor: Cadena de caracteres (20) \\
\hspace*{1cm} Mensaje: Cadena de caracteres (50) \\[12pt]


\textcolor{azul}{\textbf{RDS5.\theRDS:}} Lista de mensajes entre ambos usuarios.\\
\hspace*{1cm} Los mismos datos que en RDR5.\theRDR\\[12pt]
\stepcounter{RDR}
\stepcounter{RDS}
% ======================= %
%      ARCHIVAR CHAT      %
% ======================= %
\section{Des/archivar chat}
\textbf{\textcolor{rojo}{\underline{RF\theRF:}}}. Des/archivar chat: Permite que un usuario seleccione un chat del que quiere o no visualizar los mensajes, activando o no su visualización\\[6pt]
\stepcounter{RF}
\textbf{\textcolor{rojo}{Entrada:}} Agente externo: usuario. Acción: solicitar la des/archivación de un chat. 
Requisito de datos de entrada \textcolor{azul}{RDE5.\theRDE}. \\[6pt]

\textbf{\textcolor{rojo}{BD:}} Requisito de datos de escritura \textcolor{azul}{RDW5.\theRDW}. \\[6pt]

\textbf{\textcolor{rojo}{Salida:}} Agente externo: usuario. Acción: confirmación resultado. 
Requisito de datos de salida: ninguno \\[12pt]

\textcolor{azul}{\textbf{RDE5.\theRDE:}} Usuario que deseamos des/archivar   \\
\hspace*{1cm} Usuario: Cadena de caracteres (20) \\[12pt]

\textcolor{azul}{\textbf{RDW5.\theRDW:}} Datos del usuario que se desea archivar/desarchivar. \\
\hspace*{1cm} Usuario a archivar: Cadena de caracteres (20) \\
\hspace*{1cm} Usuario que archiva: Cadena de caracteres (20) \\
\hspace*{1cm} Usuario a archivar: Cadena de caracteres (20) \\[12pt]

\stepcounter{RDE}
\stepcounter{RDW}
% ======================== %
% RESTRICCIONES SEMÁNTICAS %
% ======================== %
\newcounter{RS} \setcounter{RS}{1}
\textbf{\textcolor{verde}{Restricciones Semánticas:}} \\[6pt]
\textbf{\textcolor{verde}{\underline{RS5.\theRS:}}} Los mensajes sólo pueden enviarse entre usuarios que son amigos. \\
\textbf{RF:} RF. \textbf{RD(s):} . \\
\hspace*{1cm} \textit{Descripción:} Si se desea enviar un mensaje a un usuario del que no se es amigo, se devolverá un error indicando que se necesita ser amigos para poder enviar un mensaje. \\[6pt]
\stepcounter{RS}

\textbf{\textcolor{verde}{\underline{RS5.\theRS:}}} Los mensajes no se pueden enviar al propio usuario. \\
\textbf{RF:} RF. \textbf{RD(s):} . \\
\hspace*{1cm} \textit{Descripción:} Si se desea enviar un mensaje al propio usuario, se devolverá un error indicando que no se puede enviar un mensaje a uno mismo. \\[6pt]
\stepcounter{RS}

\textbf{\textcolor{verde}{\underline{RS5.\theRS:}}} Para el envío de mensajes, el mensaje debe tener al menos un carácter. \\
\textbf{RF:} RF. \textbf{RD(s):} . \\
\hspace*{1cm} \textit{Descripción:} Si se desea enviar un mensaje, dicho mensaje no puede estar vacío. \\ [6pt]
\stepcounter{RS}

\textbf{\textcolor{verde}{\underline{RS5.\theRS:}}} Para el envío de mensajes, el usuario receptor debe existir. \\
\textbf{RF:} RF. \textbf{RD(s):} . \\
\hspace*{1cm} \textit{Descripción:} Si se desea enviar un mensaje, el receptor debe existir. \\ [6pt]
\stepcounter{RS}


